\chapter{Python 解析の高速化}

\section{まず測る}

高速化を始める前に、どこが遅いのかを測る。読込みが遅いのか、探索が遅いのか、図の作成が遅いのかを区別しないまま最適化しても、効果は見えにくい。

ここで役立つのが \code{time.perf\_counter} である。処理の前後で時刻を取り、差を記録するだけでも、どの段階が重いかの見当が付く。大切なのは、全体時間だけでなく、読込み、前処理、探索、書き出し、作図のように段階を分けて測ることである。

\begin{lstlisting}[language=Python, caption={処理時間を測る最小例}]
import time

t0 = time.perf_counter()
run_analysis()
t1 = time.perf_counter()
print(f"elapsed = {t1 - t0:.3f} s")
\end{lstlisting}

測定するときは、一度だけの結果を絶対視しないことも重要である。キャッシュや OS の状態で多少は揺れるため、条件をそろえ、複数回の傾向を見る方がよい。

\section{Jupyter では \code{\%timeit} も使える}

ノートブック環境で短い式や関数の速度を見たいときは、Jupyter の \code{\%timeit} や \code{\%\%timeit} も便利である。複数回の実行を自動で繰り返し、平均的な実行時間を表示してくれるため、小さな差を見たいときに役立つ。

\begin{lstlisting}[language=Python, caption={Jupyter の timeit の例}]
%%timeit
result = array * 2
\end{lstlisting}

ただし、\code{\%timeit} はノートブックや IPython での簡便な測定であり、スクリプト全体の読込みや入出力を含めた時間を見るには向かない。小さな部品には \code{\%timeit}、解析全体には \code{perf\_counter} というように使い分けると整理しやすい。

\section{計算量の見積もり}

高速化では、実測だけでなく理屈の見積もりも役に立つ。入力サイズを $N$ としたとき、処理回数が $N$ に比例するのか、$N \log N$ 程度なのか、$N^2$ なのかで、大きなデータに対する伸び方がまるで違う。

厳密な漸近解析までできなくても、「データを 2 倍にしたら計算量は何倍くらいになるか」を考えるだけで十分に有益である。高速化はテクニック集ではなく、問題の構造を考える作業である。

\begin{lstlisting}[numbers=none, xleftmargin=2\zw, caption={計算量の見方の例}]
O(N)      : 1 回走査
O(N log N): ソートを含む処理
O(N^2)    : 全組合せ比較
\end{lstlisting}

この違いは、小さなテストデータでは目立たなくても、大きな入力で急に効いてくる。だからこそ、実測の前に理屈として見積もる意味がある。

\section{遅い例と速い例を比べる}

同じ問題でも、書き方の差ではなく、アルゴリズムの差で桁違いに時間が変わることがある。この章では、その違いを具体例で見る。図~\ref{fig:vectorization-benchmark} は、単純な隣接差分フィルタを純粋な Python ループ、NumPy、pandas で処理した例である。

この点は初学者ほど見落としやすい。変数名を短くしたり、関数を減らしたりすることは、ほとんど速度に効かない。一方で、比較回数を減らす、探索範囲を絞る、配列演算へ置き換えるといった変更は、本質的に計算量を変えることがある。

\section{ペア探索の遅い例}

時刻が昇順に並んだ値が $N$ 個あるとき、すべての組 $(i, j)$ を比べる単純な二重ループは、おおよそ $N^2$ に比例する計算になる。

\begin{lstlisting}[language=Python, caption={遅いペア探索の例}]
for i in range(len(times)):
    for j in range(i + 1, len(times)):
        if abs(times[j] - times[i]) <= window:
            ...
\end{lstlisting}

この書き方は分かりやすいが、データ量が増えるとすぐに実用でなくなる。

二重ループが悪いのではなく、「不要な比較まで全部やっている」ことが問題である。小さな入力では問題なくても、データ量が 10 倍になると比較回数はおよそ 100 倍になる。この増え方を直感として持っておくと、高速化の議論が理解しやすい。

\section{ペア探索の速い例}

時刻順に並んでいる列なら、遠すぎる要素まで比べる必要はない。差が閾値を超えた時点で比較を打ち切る、あるいは隣接要素だけを見る、といった工夫で計算回数を大きく減らせる。

実際には、どこまで比較すれば十分かは問題設定に依存する。だが共通しているのは、「入力の構造を利用する」という発想である。ソート済み、単調増加、局所性がある、といった性質が分かっているなら、それを使わないのはもったいない。

\section{トリプレット探索を題材にしたアルゴリズムの差}

トリプレット探索では、単純な三重ループはさらに厳しい。ここで重要なのは、組合せを力ずくで全部試すのではなく、近傍候補を先に絞り、そこから有効な組だけを組み立てることである。

この視点は、ペア探索にも一般のクラスタ探索にも共通している。

アルゴリズムを考えるときは、まず最も素朴な実装を書いて正しさを確かめ、その後で不要な組合せを減らす方法を探すとよい。最初から複雑な高速版だけを書こうとすると、速いが間違っているコードになりやすい。

\section{NumPy と pandas による高速化}

アルゴリズムを見直した後は、NumPy や pandas へ処理を移すことで、Python レベルのループを減らせる。特に差分計算、ソート、列抽出、フィルタは、これらのライブラリが得意とする。図~\ref{fig:vectorization-benchmark} からも、配列演算が速度に与える効果を視覚的に確認できる。

ここで注意したいのは、NumPy や pandas を使えば自動的に何でも速くなるわけではないことである。根本的に無駄な比較をしているなら、配列演算へ移しても限界がある。まず計算量を減らし、そのうえで実装をベクトル化するという順序が自然である。

\section{メモリとコピーの管理}

速度だけを見ていると、メモリ使用量が大きくなりすぎて逆に遅くなることがある。巨大な中間配列を何度も作る実装では、CPU よりもメモリ転送が支配的になる場合がある。

そのため、高速化では「何回ループしたか」だけでなく、「何回コピーしたか」も見る必要がある。不要なコピーを減らし、必要な列や範囲だけを扱う設計は、速度と安定性の両方に効く。

\section{分割読込み、並列化、ストリーミング}

大規模解析では、全部を一度にメモリへ載せない設計も重要である。ファイル単位で処理し、必要な情報だけを持ち回る。さらに、独立な単位に分割できるなら並列化も検討できる。

ストリーミング処理では、「今必要な範囲だけを覚える」という考え方が中心になる。全履歴が不要なら保持しない、集計済みの量は逐次更新する、といった設計が有効である。メモリ節約と速度向上が同時に得られることも多い。

\section{プロファイラを使う}

本格的に遅い理由を調べるなら、プロファイラも役立つ。例えば \code{cProfile} を使うと、どの関数に時間が使われているかを確認できる。

\begin{lstlisting}[language=Python, caption={cProfile の最小例}]
import cProfile

cProfile.run("run_analysis()", sort="cumtime")
\end{lstlisting}

ただし、プロファイラの数字も文脈なしには読めない。回数が多いから遅いのか、1 回が重いのか、入出力を待っているのかを区別する必要がある。測定値とコード構造を一緒に見ることが重要である。

\begin{lstlisting}[numbers=none, xleftmargin=2\zw, caption={プロファイル結果の見方の例}]
ncalls  tottime  cumtime  function
100000   1.25     1.40    check_pair
    10   0.02     2.10    read_all_files
\end{lstlisting}

この結果例では、\code{check\_pair} は 1 回は軽いが回数が多く、\code{read\_all\_files} は内部の処理込みで累積時間が大きいと読める。どの列が何を意味するかを理解すると、改善方針が立てやすくなる。

\section{C や C++ へ書き換えるという選択肢}

Python は開発速度が高いが、計算コアのすべてを Python で書く必要はない。アルゴリズムを十分に整理したうえで、どうしても重い部分だけを C や C++ へ移す、というのは自然な選択肢である。

ただし、言語を変える前に Python 側で何が遅いかを把握しておく必要がある。ボトルネックが入出力や高水準の設計にあるなら、C++ 化しても大きな改善にならない。低レベル言語への移行は、最後の切り札として考える方がよい。

\begin{lstlisting}[language=C++, caption={重いループを C++ へ出す発想の例}]
for (std::size_t i = 0; i < n; ++i) {
    for (std::size_t j = i + 1; j < n; ++j) {
        if (std::abs(times[j] - times[i]) <= window) {
            ++count;
        }
    }
}
\end{lstlisting}

この書き換えは、定数倍の改善をもたらすことがある。しかし、比較回数そのものが多すぎるなら本質的な改善ではない。アルゴリズムと実装言語を分けて考える姿勢が必要である。

\section{Python から C/C++ を呼び出す}

方法はいくつかある。たとえば \code{ctypes} で共有ライブラリを呼ぶ方法、Python C 拡張を書く方法、pybind11 を使って C++ を Python から自然に使える形へ包む方法などがある。さらに、Python と C の中間に位置する Cython を使い、Python に近い構文のまま重いループだけをコンパイルする方法も、この節に含まれる重要な選択肢である。

\begin{lstlisting}[language=Python, caption={Cython で型を付ける発想の例}]
cdef int i
cdef double total = 0.0

for i in range(n):
    total += values[i]
\end{lstlisting}

どの方法を選ぶかは、必要な速度、保守性、ビルドの難しさで決まる。短い試作なら \code{ctypes}、Python らしい使い勝手まで整えたいなら pybind11、既存の Python 風コードを保ちながら内側だけ速くしたいなら Cython、といった選択が考えられる。重要なのは、インタフェースを狭く保ち、重い部分だけを外へ出すことである。

\begin{lstlisting}[language=Python, caption={Python 側の呼び出しイメージ}]
result = fast_module.count_pairs(times, window)
\end{lstlisting}

Python 側の呼び出しが単純であれば、内部実装を後から差し替えやすい。まず Python で正しい API を決め、その内部だけを高速化するという順が自然である。なお、Cython も強力だが、アルゴリズムそのものが無駄な場合には決定打にならない。この点は C++ 化と同じであり、まず Python 側で何が遅いのか、どのループが本当に支配的なのかを見極めた後に使う方が効果的である。

\section{外部実装を参考にするときの見方}

高速化を学ぶときは、外部の C や C++ の実装例を読むことが参考になる。ただし、そこで大切なのは固有名詞を覚えることではない。まず見るべきなのは、どの問題を解いているのか、どのアルゴリズムを選んでいるのか、どの部分だけを低レベル言語へ移しているのか、という点である。

高速化の章では、このような例を参照しながら、次の二つを常に区別して考える。

\begin{itemize}
\item アルゴリズムを変えたから速くなったのか
\item 実装言語を変えたから速くなったのか
\end{itemize}

多くの場合、最初に効くのは前者である。

したがって、高速化のレポートを書くときは、「NumPy にした」「C++ にした」で終わらせない方がよい。何が変わった結果として速くなったのかを、計算量、比較回数、メモリ配置、ループの位置といった観点で説明するべきである。

\section{結果例と報告の仕方}

高速化の結果は、本文中で比較できる形に整理すると読みやすい。例えば、次のような簡潔な結果でも十分に意味がある。

\begin{lstlisting}[numbers=none, xleftmargin=2\zw, caption={速度比較の結果例}]
slow version : 182.4 s
vectorized   :  12.7 s
streaming    :   4.8 s
\end{lstlisting}

この表現だけでは理由が分からないので、続けて「比較回数を減らした」「Python ループを配列演算へ移した」「全件保持をやめた」といった説明を添える。数値と理由の両方がそろって初めて報告になる。

\begin{lstlisting}[language=Python, caption={隣接差分を NumPy でまとめて計算する例}]
import numpy as np

values = np.array(data, dtype=np.float64)
diffs = np.diff(values)
count = np.count_nonzero(diffs > 0.25)
\end{lstlisting}

\begin{figure}[tbp]
\centering
\includegraphics[width=0.82\linewidth]{vectorization_benchmark.pdf}
\caption{同じ隣接差分フィルタを、純粋な Python ループ、NumPy、pandas で実行した例。縦軸は実行時間で、対数目盛を使っている。}
\label{fig:vectorization-benchmark}
\end{figure}
